% GENERATED FILE. Edit canonical lessons, then run npm run generate:books.
% canonical-source-hash: fnv1a64:b6dc266faf40480b
% canonical-lessons: TA-C322-manamuti, TA-C322-tali, TA-C322-mukurttam, TA-C322-manametai, TA-C322-pantal

\chapter{A Wedding Day}
\label{ch:ta-a-wedding-day}
\hlchaptermodality{322}

\begin{chapteropening}
\textbf{By the end of this chapter:} \emph{I can describe a wedding day.}

Complete the last lesson of chapter 322: I can describe a wedding day.
\end{chapteropening}

\section[\texorpdfstring{\ta{மணமுடி} (maṇamuṭi)}{மணமுடி (maṇamuṭi)}]{\ta{மணமுடி} (maṇamuṭi) — to wed; to marry}

\label{lesson:TA-C322-manamuti}



\begin{quote}
Before the new one: say the Tamil for a baby shower, then the Tamil for a housewarming.
\end{quote}

\subsection*{You'll want to know: \ta{மணமுடி}}
\textbf{\ta{மணமுடி}} — \emph{maṇamuṭi} — \textquotedblleft{}to wed; to marry\textquotedblright{}.

\begin{quote}
\ta{அவர்} \ta{கடந்த} \ta{ஆண்டு} \ta{ஒரு} \ta{ஆசிரியையை} \ta{மணமுடித்தார்}.
\end{quote}

\begin{quote}
\emph{avar kaṭanta āṇṭu oru āciriyaiyai maṇamuṭittār.}
\end{quote}

\begin{quote}
Last year he married a teacher.
\end{quote}

Past (I …): \ta{மணமுடித்தேன்}, \emph{maṇamuṭittēṉ}.

A written word for \ta{திருமணம்} \ta{செய்}; the spouse takes -\ta{ஐ}.

\begin{cousinweb}
\textbf{\ta{மணமுடி}} (\emph{maṇamuṭi}) is a Tamil compound built from \ta{மணம்}, \textquotedblleft{}union, marriage\textquotedblright{} (DEDR 4667), and \ta{முடி}, \textquotedblleft{}to tie, fasten\textquotedblright{} (DEDR 4921).
\end{cousinweb}

\begin{grammarlens}[title={What you've built}]
One more word for telling what happened.
\end{grammarlens}

\subsection*{Your turn}
\begin{itemize}
\raggedright
  \item \emph{Say it:} \emph{maṇamuṭi}
  \item \emph{Say it:} \emph{maṇamuṭi}, once more
  \item \emph{Say it:} \emph{avar kaṭanta āṇṭu oru āciriyaiyai maṇamuṭittār.} — Last year he married a teacher.
  \item \emph{Recall:} say the Tamil for a baby shower, then the Tamil for a housewarming, then say \emph{maṇamuṭi} again
\end{itemize}

\begin{tcolorbox}[breakable,colback=teal!4,colframe=teal!35!black,title={Before you move on}]
What does \ta{மணமுடி} mean? (\textquotedblleft{}To wed; to marry\textquotedblright{}.) Say it once more.
\end{tcolorbox}
\section[\texorpdfstring{\ta{தாலி} (tāli)}{தாலி (tāli)}]{\ta{தாலி} (tāli) — a wedding pendant}

\label{lesson:TA-C322-tali}



\begin{quote}
Before the new one: say the Tamil for a housewarming, then the Tamil for to wed.
\end{quote}

\subsection*{You'll want to know: \ta{தாலி}}
\textbf{\ta{தாலி}} — \emph{tāli} — \textquotedblleft{}a wedding pendant\textquotedblright{}.

\begin{quote}
\ta{மணமகன்} \ta{மணமகளுக்குத்} \ta{தாலி} \ta{கட்டினார்}.
\end{quote}

\begin{quote}
\emph{maṇamakaṉ maṇamakaḷukkut tāli kaṭṭiṉār.}
\end{quote}

\begin{quote}
The groom tied the wedding pendant on the bride.
\end{quote}

\ta{தாலி} \ta{கட்டு}, to tie the thali, is the moment of marriage.

\begin{cousinweb}
\textbf{\ta{தாலி}} (\emph{tāli}) is native Dravidian (DEDR 3175): the neck ornament the groom ties on the bride.
\end{cousinweb}

\begin{grammarlens}[title={What you've built}]
One more word for telling what happened.
\end{grammarlens}

\subsection*{Your turn}
\begin{itemize}
\raggedright
  \item \emph{Say it:} \emph{tāli}
  \item \emph{Say it:} \emph{tāli}, once more
  \item \emph{Say it:} \emph{maṇamakaṉ maṇamakaḷukkut tāli kaṭṭiṉār.} — The groom tied the wedding pendant on the bride.
  \item \emph{Recall:} say the Tamil for a housewarming, then the Tamil for to wed, then say \emph{tāli} again
\end{itemize}

\begin{tcolorbox}[breakable,colback=teal!4,colframe=teal!35!black,title={Before you move on}]
What does \ta{தாலி} mean? (\textquotedblleft{}A wedding pendant\textquotedblright{}.) Say it once more.
\end{tcolorbox}
\section[\texorpdfstring{\ta{முகூர்த்தம்} (mukūrttam)}{முகூர்த்தம் (mukūrttam)}]{\ta{முகூர்த்தம்} (mukūrttam) — an auspicious hour}

\label{lesson:TA-C322-mukurttam}



\begin{quote}
Before the new one: say the Tamil for to wed, then the Tamil for a wedding pendant.
\end{quote}

\subsection*{You'll want to know: \ta{முகூர்த்தம்}}
\textbf{\ta{முகூர்த்தம்}} — \emph{mukūrttam} — \textquotedblleft{}an auspicious hour\textquotedblright{}.

\begin{quote}
\ta{முகூர்த்தம்} \ta{காலை} \ta{ஒன்பது} \ta{மணிக்கு}.
\end{quote}

\begin{quote}
\emph{mukūrttam kālai oṉpatu maṇikku.}
\end{quote}

\begin{quote}
The auspicious hour is at nine in the morning.
\end{quote}

Also the wedding itself: \ta{முகூர்த்த} \ta{நாள்}, the wedding day.

\begin{cousinweb}
\textbf{\ta{முகூர்த்தம்}} (\emph{mukūrttam}) is borrowed from Sanskrit \emph{muhūrta}, \textquotedblleft{}a moment, a fixed period\textquotedblright{} (CDIAL 10229).
\end{cousinweb}

\begin{grammarlens}[title={What you've built}]
One more word for telling what happened.
\end{grammarlens}

\subsection*{Your turn}
\begin{itemize}
\raggedright
  \item \emph{Say it:} \emph{mukūrttam}
  \item \emph{Say it:} \emph{mukūrttam}, once more
  \item \emph{Say it:} \emph{mukūrttam kālai oṉpatu maṇikku.} — The auspicious hour is at nine in the morning.
  \item \emph{Recall:} say the Tamil for to wed, then the Tamil for a wedding pendant, then say \emph{mukūrttam} again
\end{itemize}

\begin{tcolorbox}[breakable,colback=teal!4,colframe=teal!35!black,title={Before you move on}]
What does \ta{முகூர்த்தம்} mean? (\textquotedblleft{}An auspicious hour\textquotedblright{}.) Say it once more.
\end{tcolorbox}
\section[\texorpdfstring{\ta{மணமேடை} (maṇamēṭai)}{மணமேடை (maṇamēṭai)}]{\ta{மணமேடை} (maṇamēṭai) — a wedding dais}

\label{lesson:TA-C322-manametai}



\begin{quote}
Before the new one: say the Tamil for a wedding pendant, then the Tamil for an auspicious hour.
\end{quote}

\subsection*{You'll want to know: \ta{மணமேடை}}
\textbf{\ta{மணமேடை}} — \emph{maṇamēṭai} — \textquotedblleft{}a wedding dais\textquotedblright{}.

\begin{quote}
\ta{மணமக்கள்} \ta{மணமேடையில்} \ta{அமர்ந்தார்கள்}.
\end{quote}

\begin{quote}
\emph{maṇamakkaḷ maṇamēṭaiyil amarntārkaḷ.}
\end{quote}

\begin{quote}
The bride and groom sat on the wedding dais.
\end{quote}

\ta{மணம்}, marriage, as in \ta{மணமகன்}, the groom.

\begin{cousinweb}
\textbf{\ta{மணமேடை}} (\emph{maṇamēṭai}) is a Tamil compound built from \ta{மணம்}, \textquotedblleft{}marriage\textquotedblright{} (DEDR 4667), and \ta{மேடை}, \textquotedblleft{}a raised platform\textquotedblright{}, whose origin is not settled in the sources at hand.
\end{cousinweb}

\begin{grammarlens}[title={What you've built}]
One more word for telling what happened.
\end{grammarlens}

\subsection*{Your turn}
\begin{itemize}
\raggedright
  \item \emph{Say it:} \emph{maṇamēṭai}
  \item \emph{Say it:} \emph{maṇamēṭai}, once more
  \item \emph{Say it:} \emph{maṇamakkaḷ maṇamēṭaiyil amarntārkaḷ.} — The bride and groom sat on the wedding dais.
  \item \emph{Recall:} say the Tamil for a wedding pendant, then the Tamil for an auspicious hour, then say \emph{maṇamēṭai} again
\end{itemize}

\begin{tcolorbox}[breakable,colback=teal!4,colframe=teal!35!black,title={Before you move on}]
What does \ta{மணமேடை} mean? (\textquotedblleft{}A wedding dais\textquotedblright{}.) Say it once more.
\end{tcolorbox}
\section[\texorpdfstring{\ta{பந்தல்} (pantal)}{பந்தல் (pantal)}]{\ta{பந்தல்} (pantal) — a pandal; a canopy}

\label{lesson:TA-C322-pantal}



\begin{quote}
Before the new one: say the Tamil for an auspicious hour, then the Tamil for a wedding dais.
\end{quote}

\subsection*{You'll want to know: \ta{பந்தல்}}
\textbf{\ta{பந்தல்}} — \emph{pantal} — \textquotedblleft{}a pandal; a canopy\textquotedblright{}.

\begin{quote}
\ta{வீட்டின்} \ta{முன்} \ta{பெரிய} \ta{பந்தல்} \ta{போட்டார்கள்}.
\end{quote}

\begin{quote}
\emph{vīṭṭiṉ muṉ periya pantal pōṭṭārkaḷ.}
\end{quote}

\begin{quote}
They put up a big pandal in front of the house.
\end{quote}

\ta{பந்தல்} \ta{போடு}, to put up a pandal; English pandal is usually traced to this word.

\begin{cousinweb}
\textbf{\ta{பந்தல்}} (\emph{pantal}) is native Dravidian (DEDR 3922): a pandal, a shed or an arbour.
\end{cousinweb}

\begin{grammarlens}[title={What you've built}]
One more word for telling what happened.
\end{grammarlens}

\subsection*{Your turn}
\begin{itemize}
\raggedright
  \item \emph{Say it:} \emph{pantal}
  \item \emph{Say it:} \emph{pantal}, once more
  \item \emph{Say it:} \emph{vīṭṭiṉ muṉ periya pantal pōṭṭārkaḷ.} — They put up a big pandal in front of the house.
  \item \emph{Recall:} say the Tamil for an auspicious hour, then the Tamil for a wedding dais, then say \emph{pantal} again
\end{itemize}

\begin{tcolorbox}[breakable,colback=teal!4,colframe=teal!35!black,title={Before you move on}]
What does \ta{பந்தல்} mean? (\textquotedblleft{}A pandal; a canopy\textquotedblright{}.) Say it once more.
\end{tcolorbox}
